% photoelectric_history.tex
% Comprehensive historical data on the photoelectric effect, 1887–1916.
% Sources: primary papers accessed via Archive.org, Royal Society, APS;
%          secondary analysis via Holton (1999) and Wikipedia reference lists.
% All numerical data verified against Millikan (1916), Phys.\ Rev.\ 7:355--388.

\documentclass[12pt]{article}
\usepackage{amsmath,amssymb}
\usepackage{booktabs}
\usepackage{array}
\usepackage{hyperref}
\usepackage[margin=1in]{geometry}
\usepackage{setspace}
\onehalfspacing

\title{Historical Data on the Photoelectric Effect:\\
Lenard (1902) through Millikan (1916)}
\author{Background notes for T.\ J.\ Sargent}
\date{\today}

\begin{document}
\maketitle
\tableofcontents
\bigskip

\section{Overview}

The photoelectric effect---the emission of electrons from a metal surface
irradiated by light---was discovered by Heinrich Hertz in 1887.  The subsequent
two decades of experimental work culminated in Einstein's 1905 quantum
hypothesis and its direct verification by Millikan in 1916.  The central
observable is the \emph{stopping potential} $V_0$: the minimum retarding voltage
needed to prevent the most energetic emitted electrons from reaching the
collector.  Einstein's equation predicts
\begin{equation}\label{eq:einstein}
  V_0 = \frac{h}{e}\,\nu - \frac{W}{e},
\end{equation}
where $h$ is Planck's constant, $e$ the elementary charge, $\nu$ the frequency
of the incident light, and $W$ the work function of the metal surface.
Equation~\eqref{eq:einstein} implies a \emph{linear relationship} between $V_0$
and $\nu$ with slope $h/e$ and (negative) intercept $-W/e$.

\section{Pre-Einstein Experimental History}

\subsection{Hertz (1887) and Hallwachs (1888)}

Heinrich Hertz observed in 1887 that ultraviolet light falling on a metal spark
gap reduced the voltage needed to produce a spark.  Wilhelm Hallwachs (1888)
showed that a zinc plate illuminated by ultraviolet light discharged negative
charge and acquired positive charge---i.e., light caused emission of negative
particles.

\subsection{Lenard's Experiments (1899--1902)}

Philipp Lenard established several qualitative properties of photoelectron
emission:
\begin{enumerate}
\item (1899) The emitted particles are identical to cathode-ray electrons,
  confirming J.\ J.\ Thomson's identification of the electron.
\item (1902) \textbf{Key result}: the \emph{maximum velocity} (equivalently,
  the stopping potential $V_0$) of emitted electrons is \emph{independent of
  light intensity} but depends on the frequency of the incident light.  This
  directly contradicted classical wave theory.
\item Lenard noted that the ``trigger'' action of light seemed to release energy
  stored in the metal rather than delivering energy directly from the wave.
\end{enumerate}

\textbf{Quantitative limitations.}  Lenard's measurements of $V_0$ vs.\ $\nu$
were qualitative rather than quantitative.  His metal surfaces were typically
oxidized, introducing variable and uncontrolled contact potentials that shifted
the apparent stopping voltages.  He did not work in ultra-high vacuum.
Consequently, while his \emph{qualitative} finding (that $V_0$ depends on
frequency, not intensity) was correct and epochal, his data could not
reliably test the linear prediction~\eqref{eq:einstein}.  The quantitative
confirmation awaited Millikan's vacuum technique developed 1908--1916.

\textbf{Reference.}  Lenard's 1902 paper: ``Ueber die lichtelektrische
Wirkung,'' \emph{Ann.\ Phys.}\ \textbf{313}(5):149--198 (1902).
\doi{10.1002/andp.19023130510}  (Series numbering: \emph{Ann.\ Phys.}
\textbf{8}:149--198 in the ``neue Folge'' [new series].)

\subsection{Other Pre-1905 Experimenters}

\paragraph{Elster and Geitel (1889--1891).}
Johann Elster and Hans Geitel demonstrated that the alkali metals (Na, K, Rb,
Cs) are photoelectrically active at visible wavelengths, whereas most metals
require ultraviolet light.  They concluded (1891) that more electropositive
metals have a smaller threshold frequency $\nu_0$ (longer threshold
wavelength $\lambda_0$), anticipating point~4 of Einstein's equation.
See: Wied.\ Ann.\ \textbf{43}:225 (1891).

\paragraph{Mikhail Stoletov (1888--1891).}
Aleksandr Stoletov made systematic measurements of photocurrent as a function
of light intensity and polarity, establishing the photocurrent--intensity
proportionality and verifying that photoelectrons carry negative charge.
His apparatus closely resembles the modern photoelectric cell.

\paragraph{J.\ J.\ Thomson (1899).}
Thomson measured the charge-to-mass ratio $e/m$ of photoelectrons and found
it identical to that of cathode-ray electrons, confirming that photoelectrons
are electrons.

\section{Einstein's 1905 Quantum Hypothesis}

In his March 1905 paper Einstein proposed that electromagnetic radiation of
frequency $\nu$ consists of discrete quanta (``light quanta,'' later
``photons'') each carrying energy $h\nu$.  A photoelectron absorbs one quantum
and escapes the metal with kinetic energy
\[
  \tfrac{1}{2}m_e v^2_{\max} = eV_0 = h\nu - W,
\]
which is exactly equation~\eqref{eq:einstein}.  At the time, ``there were
available no experiments whatever for determining anything about how P.D.\ varies
with $\nu$'' (Millikan, 1916, p.~356).  Einstein's prediction was considered
``bold, not to say the reckless, hypothesis'' (Millikan's phrase, echoing the
general consensus).

\textbf{Reference.}  Einstein (1905): ``\"Uber einen die Erzeugung und
Verwandlung des Lichtes betreffenden heuristischen Gesichtspunkt,''
\emph{Ann.\ Phys.}\ \textbf{17}:132--148 (1905).  \doi{10.1002/andp.19053220607}

\section{Pre-Millikan Tests (1907--1914)}

Several experiments tested aspects of Einstein's equation but fell short of
definitive confirmation, primarily because of two systematic errors: (a) stray
scattered light of shorter wavelength contaminating the stopping-potential
determination, and (b) unknown contact potentials between the photoemitting
surface and the collector (which shift the intercept $-W/e$ by a constant but
do not affect the \emph{slope} $h/e$ unless they vary between measurements).

\paragraph{Ladenburg (1907).}  First quantitative measurements of $V_0$
vs.\ $\nu$, but over a wavelength range of only 50~nm (2000--2600\,\AA),
making it impossible to discriminate $V_0 \propto \nu$ (Einstein) from
$V_0 \propto \nu^{1/2}$ (classical).

\paragraph{Hughes (1912).}  A.\ Ll.\ Hughes measured stopping potentials at
three frequencies (corresponding to $\lambda = 1849$, 2257, and 2537\,\AA)
for several metals using a refined vacuum apparatus.  Two of the three points
were used to fix the regression line; the third lay close to it.  Using a
two-frequency fit, Hughes found slope values ranging from $3.55\times10^{-27}$
to $5.85\times10^{-27}$ erg$\cdot$s across eight metals (compared to
Planck's $h = 6.55\times10^{-27}$ erg$\cdot$s), a discrepancy of 10--45\%.
Systematic errors (contact potentials unmeasured in vacuo; only three
frequencies covering a 60~nm range; scattered short-wave-length light)
prevented a precise determination.

\textbf{Reference.}  Hughes (1913): ``On the emission velocities of
photo-electrons,'' \emph{Phil.\ Trans.\ R.\ Soc.\ A}
\textbf{212}:205--226 (1913). \doi{10.1098/rsta.1913.0007}

\paragraph{Richardson and Compton (1912).}  Owen Richardson and Karl Compton
made measurements on eight metals.  Their slopes (determining $h$) varied by
more than 60\% across metals, which Millikan attributed chiefly to the
restricted wavelength range (all their curves save Na and Al were between
2100 and 2750\,\AA, just 65~nm wide) and to lack of contact-EMF correction
for sodium.  Their sodium result gave $h \approx 5.2\times10^{-27}$ erg$\cdot$s
(20\% low) and their aluminum result $h \approx 4.3\times10^{-27}$ erg$\cdot$s
(35\% low).

\textbf{Reference.}  Richardson and Compton (1912): ``The photoelectric
effect,'' \emph{Phil.\ Mag.}\ \textbf{24}:575--594 (1912).
\doi{10.1080/14786441208637039}

\paragraph{Millikan (1914, preliminary).}  Millikan reported preliminary
results to the American Physical Society in April 1914 (sodium data only).
\textbf{Reference.}  Millikan (1914): ``A Direct Determination of `h',''
\emph{Phys.\ Rev.}\ \textbf{4}(1):73--75 (1914).
\doi{10.1103/PhysRev.4.73.2}

\section{Millikan's Definitive Measurement (1916)}

\subsection{Apparatus: the ``Machine Shop in Vacuo''}

Millikan's key innovation was eliminating surface contamination by building
what he called a ``machine shop in vacuo.''  Inside an evacuated bulb, an
electromagnetically controlled rotating knife (K) was used to cut a fresh
shaving from a cast cylinder of sodium, potassium, or lithium mounted on a
wheel (W).  The wheel was then rotated to bring the freshly scraped surface
into alignment with the monochromatic light beam.  All metal surfaces, knife,
and wheel were controlled entirely by external electromagnets---no mechanical
feedthroughs that might cause vacuum leaks.  This was the ninth tube of the
general design; the mechanic Julius Pearson is explicitly credited.

The alkali metals studied were \textbf{sodium} (Na) and \textbf{lithium} (Li).
(Potassium data were lost to an accident before publication.)

The light source was a Heracus 200-volt quartz-mercury arc.  A Hilger
monochromator selected individual spectral lines.  Only spectral lines that
had no companion lines on the short-wave-length side (i.e., no contamination
by higher-frequency light) were used.

Stopping potentials were measured by the Kelvin method (Kelvin double-balance):
the potentiometer voltage between the alkali metal and the collecting plate S
was adjusted until lifting the plate produced no deflection in a connected
electrometer.  The electrometer sensitivity was 2500 mm/V; the Weston Standard
Laboratory voltmeter used was certified by the National Bureau of Standards to
0.1\%.

\subsection{Mercury Arc Lines Selected}

The following mercury spectral lines were used (wavelengths determined by
Walter T.\ Whitney using a grating photograph):

\begin{table}[ht]
\centering
\caption{Mercury arc spectral lines used by Millikan (1916).
Frequencies computed from $\nu = c/\lambda$ with
$c = 2.998\times10^{10}$\,cm/s.}
\label{tab:hg_lines}
\begin{tabular}{ccc}
\toprule
$\lambda$ (\AA) & $\lambda$ (nm) & $\nu$ ($10^{14}$\,Hz) \\
\midrule
5461.0 & 546.10 & 5.490 \\
4339.4 & 433.94 & 6.908 \\
4046.8 & 404.68 & 7.409 \\
3650.2 & 365.02 & 8.213 \\
3125.5 & 312.55 & 9.591 \\
2534.7 & 253.47 & 11.827 \\
2399   & 239.9  & 12.497 \\
\bottomrule
\end{tabular}
\end{table}

\subsection{Systematic Error Controls}

\paragraph{Stray short-wave-length light.}  For each spectral line, filters
were interposed to absorb all light of shorter wavelength.  Without this
precaution, faint scattered UV light from the monochromator prism and lenses
would produce photoelectrons at all wavelengths, pushing the apparent
stopping potential for longer (lower-energy) lines upward and flattening the
slope \emph{below} $h/e$.  Millikan showed this was a primary reason previous
workers (Ladenburg, Richardson \& Compton, Hughes) obtained values of $h$
10--45\% below Planck's value.

\paragraph{Reverse electronic current.}  If the wavelength used \emph{exceeds}
the threshold wavelength of the collecting cylinder, electrons liberated from
the cylinder walls by reflected light create a ``back leak'' of current in the
wrong direction, pushing the apparent intercept toward shorter wavelengths.
Millikan used copper-oxide walls whose long-wavelength limit was
$\lambda = 2688$\,\AA\ (for the sodium tube) and $\lambda \approx 2535$\,\AA\
(for the lithium tube), and corrected the line 2535 reading accordingly.
Lines 3125 through 5461, which lie above the oxide threshold, were free of
this error.

\paragraph{Contact electromotive force.}  The contact EMF between alkali
metal and collector varied from 0.4 to 3.0\,V depending on surface age and
absorbed gas, and changed rapidly in the minutes after shaving.  The contact
EMF acts as an additional retardation or acceleration for photoelectrons.
With all lines held constant in every other respect, the contact EMF shifts the
entire $V_0$--$\nu$ curve vertically but does \emph{not} change its slope.
Thus the slope $h/e$ can be determined accurately even if the contact EMF is
not independently measured---provided the EMF is \emph{constant} during the
full series of measurements.

Millikan found that with $\sim 0.01$~mm Hg of residual gas in the bulb,
the contact EMF stabilised within about one hour of shaving and remained
constant to within $\pm 0.01$~V for several hours---ample time for a
complete set of measurements.  In the main sodium run, the contact EMF
was measured independently by the Kelvin method as $V_c = 2.51$\,V
(Na positive with respect to Cu oxide), remaining stable from
2.495\,V (6\,h after shaving) to 2.493\,V (9\,h after shaving).

\subsection{Raw Data: Sodium (Table I)}

Millikan tabulated the electrometer deflections (in mm, for a 30-second
illumination) as a function of applied potentiometer voltage for each
spectral line (Table I of the original paper).  The balance point
(zero deflection) gives the intercept.  Deflections at the last few settings
before the balance for the main sodium curves are reproduced below.

\begin{table}[ht]
\centering
\caption{Selected readings from Millikan (1916), Table~I.  Sodium surface.
``Volts'' is the applied potentiometer voltage; ``Def'n'' is the
electrometer deflection in mm (30-s illumination).  For lines 5461--3126
the contact EMF ($V_c = 2.51$\,V) exceeds the photopotential so the
balance falls on the negative-voltage side; the magnitudes of the applied
\emph{assisting} potential near balance are listed.  For line 2535 the
balance falls on the positive (retarding) side.}
\label{tab:millikan_raw}
\begin{tabular}{cc|cc|cc|cc|cc|cc}
\toprule
\multicolumn{2}{c|}{5461\,\AA} &
\multicolumn{2}{c|}{4339\,\AA} &
\multicolumn{2}{c|}{4047\,\AA} &
\multicolumn{2}{c|}{3650\,\AA} &
\multicolumn{2}{c|}{3126\,\AA} &
\multicolumn{2}{c}{2535\,\AA} \\
V (V) & mm & V (V) & mm & V (V) & mm & V (V) & mm & V (V) & mm & V (V) & mm \\
\midrule
2.257 & 28 & 1.581 & 44 & 1.576 & 82 & 1.157 & 67.5 & 0.581 & 52 & $-$0.058 & 68 \\
2.205 & 14 & 1.629 & 20 & 1.524 & 55 & 1.105 & 36   & 0.529 & 29 & $+$0.058 & 38 \\
2.152 &  7 & 1.576 & 10 & 1.471 & 36 & 1.053 & 19   & 0.477 & 12 & $+$0.162 & 26 \\
2.100 &  3 & 1.524 &  4 & 1.419 & 24 & 1.000 & 11   & 0.424 &  5 & $+$0.267 & 16.5\\
      &   & 1.367 &  3 & 0.948 &  4 &       &      & 0.372 & 2.5& $+$0.372 &  8 \\
\bottomrule
\end{tabular}
\end{table}

\subsection{Determination of \texorpdfstring{$h$}{h} from Sodium}

\paragraph{Main five-point fit (Fig.\ 6 of the paper).}
The balance-point intercepts from the curves of Fig.\ 5 were plotted
against frequency $\nu$.  The slope of the resulting straight line,
determined primarily from the five most reliable points
($\lambda = 5461, 4339, 4047, 3650, 3126$\,\AA), was
\[
  \left.\frac{V_0}{\nu}\right|_{\text{slope}} =
  4.124 \times 10^{-15}\text{ V$\cdot$s}
  \quad \text{(Fig.\ 6 fit, sodium, Millikan 1916, p.\ 374).}
\]
``No potential departs from the line by more than 0.01~V and since the
range of volts is about 2, it is a conservative estimate to place the
maximum uncertainty in the slope at about 1 in 200 or 0.5 per cent''
(Millikan 1916, p.\ 374).

\paragraph{Nine two-point determinations (Table II of the paper).}
To check for consistency and eliminate the influence of a slowly drifting
contact EMF, Millikan also measured the slope from pairs of lines on
different days.  Line 5461 was always one of the two (chosen for its
large intensity and hence most precise intercept).  The nine results are
reproduced in Table~\ref{tab:slope_determinations}.

\begin{table}[ht]
\centering
\caption{Millikan (1916): Table~II.  Nine independent two-point
determinations of the slope $h/e$ from sodium.  Each row is from
a different day and a different vacuum condition.
The ``estimated uncertainties'' quoted by Millikan at the time of
measurement range from 1\% to 4\%, with lines 3650 and 3126 most
reliable.}
\label{tab:slope_determinations}
\begin{tabular}{cc}
\toprule
$\lambda$ compared with 5461\,\AA & Slope ($\times 10^{-15}$\,V$\cdot$s) \\
\midrule
3126  & 4.11 \\
3650  & 4.14 \\
3126  & 4.10 \\
3650  & 4.12 \\
3126  & 4.24 \\
4047  & 3.98 \\
2535  & 4.04 \\
3126  & 4.24 \\
4047  & 4.21 \\
\midrule
Mean  & $\mathbf{4.131}$ \\
\bottomrule
\end{tabular}
\end{table}

\paragraph{Combined sodium result.}
Taking the mean of 4.124 (from the five-point fit) and 4.131 (from
Table~\ref{tab:slope_determinations}), Millikan obtained the combined
sodium slope:
\[
  \frac{h}{e} = 4.128 \times 10^{-15}\text{ V$\cdot$s}.
\]
With his oil-drop value $e = 4.774\times10^{-10}$ esu~$= 1.592\times10^{-19}$\,C
this gives
\[
  h_{\text{Na}} = 4.128\times10^{-15}\times 1.592\times10^{-19}
  = 6.572\times10^{-34}\text{ J$\cdot$s}
  \approx 6.57\times10^{-27}\text{ erg$\cdot$s.}
\]

\subsection{Lithium Results}

For lithium the measurements were easier: the photosensitivity remained
high for months without reshaving.  Two independent series gave:

\paragraph{First lithium series (Table III of the paper).}
Contact EMF = 1.52\,V; threshold frequency $\nu_0(\text{Li}) = 57.0\times10^{13}$\,Hz
($\lambda_0 \approx 5260$\,\AA).  At $\lambda = 2535$\,\AA: $V_{0,\text{obs}} = 1.00$\,V,
so $V_{0,\text{true}} = 1.00 + 1.52 = 2.52$\,V.

\paragraph{Second lithium series (six months later, Fig.\ 9 of the paper).}
Contact EMF = 1.11\,V; $\nu_0(\text{Li}) = 59.7\times10^{13}$\,Hz
($\lambda_0 \approx 5020$\,\AA).  At $\lambda = 2535$\,\AA: $V_{0,\text{obs}} = 1.29$\,V,
so $V_{0,\text{true}} = 1.29 + 1.11 = 2.40$\,V.

From the two lithium series, Millikan obtained:
\[
  h_{\text{Li}} = 6.584\times10^{-27}\text{ erg$\cdot$s}
  \quad
  (\text{mean of two determinations; uncertainty}\approx1\%).
\]

\subsection{Final Result}

Placing greater reliance on the sodium determination (smaller uncertainty):
\begin{equation}\label{eq:h_result}
  \boxed{h = 6.57\times10^{-27}\text{ erg$\cdot$s}
        = 6.57\times10^{-34}\text{ J$\cdot$s}}
  \quad
  (\text{precision } \approx 0.5\%).
\end{equation}
The modern (NIST CODATA 2022) value is
$h = 6.626\,070\,15\times10^{-34}$\,J$\cdot$Hz$^{-1}$ (exact by SI definition
since 2019).  Millikan's result is low by about 0.8\%---well within his
stated uncertainty---attributable primarily to the small error in his
oil-drop value of $e$ (Millikan used $e = 1.592\times10^{-19}$\,C; the modern
value is $e = 1.602\,176\,634\times10^{-19}$\,C, about 0.65\% higher).

\section{Reconstructed \texorpdfstring{$(\nu, V_0)$}{(nu, V0)}
         Data Table for Linear Regression}

The \emph{true} stopping potentials in Einstein's equation~\eqref{eq:einstein}
are $V_{0,\text{true}} = V_{0,\text{obs}} + V_c$, where $V_c$ is the contact
EMF, which shifts the curve vertically but does not affect the slope.
These true stopping potentials can be reconstructed from the known
slope $h/e = 4.128\times10^{-15}$\,V$\cdot$s, the independently
determined threshold frequency $\nu_0$, and the explicitly stated
$V_{0,\text{true}}$ for the line 2535\,\AA.

The relevant parameters are:
\begin{align*}
  \nu_0(\text{Na}) &= 4.39\times10^{14}\text{ Hz}
    \quad (\lambda_0 = 6800\text{\,\AA, from contact-EMF analysis}), \\
  \nu_0(\text{Li}) &= 5.70\times10^{14}\text{ Hz}
    \quad (\lambda_0 = 5260\text{\,\AA, first series}), \\
  W_{\text{Na}} &= h\nu_0 = 6.57\times10^{-27}\times4.39\times10^{13}
                 = 2.88\times10^{-12}\text{ erg} \approx 1.80\text{ eV}, \\
  W_{\text{Li}} &= 6.57\times10^{-27}\times5.70\times10^{13}
                 = 3.74\times10^{-12}\text{ erg} \approx 2.34\text{ eV}.
\end{align*}

\begin{table}[ht]
\centering
\caption{Reconstructed true stopping potentials $V_0$ for sodium and lithium,
computed from $V_0 = (h/e)(\nu - \nu_0)$ with $h/e = 4.128\times10^{-15}$\,V$\cdot$s.
For sodium: $\nu_0 = 4.39\times10^{14}$\,Hz.  For lithium: $\nu_0 = 5.70\times10^{14}$\,Hz
(first series).  The explicit data point for Na at 2535\,\AA\ is stated in
the paper as $V_{0,\text{obs}} = 0.52$\,V and $V_c = 2.51$\,V, giving
$V_{0,\text{true}} = 3.03$\,V (computed: 3.07\,V; the small difference reflects
scattered-light correction for this line noted by Millikan).}
\label{tab:stopping_potentials}
\begin{tabular}{ccccc}
\toprule
$\lambda$ (\AA) & $\lambda$ (nm) & $\nu$ ($10^{14}$\,Hz)
  & $V_0$(Na) (V) & $V_0$(Li) (V) \\
\midrule
5461.0 & 546.10 & 5.490 & 0.454 & ---$^{a}$ \\
4339.4 & 433.94 & 6.908 & 1.039 & 0.499 \\
4046.8 & 404.68 & 7.409 & 1.246 & 0.706 \\
3650.2 & 365.02 & 8.213 & 1.578 & 1.037 \\
3125.5 & 312.55 & 9.591 & 2.147 & 1.606 \\
2534.7 & 253.47 & 11.827 & $3.07^{b}$ & 2.529 \\
\bottomrule
\end{tabular}
\begin{tablenotes}
  \footnotesize
  \item[$a$] Line 5461 lies above the long-wavelength limit of the
  lithium surface (Li threshold $\lambda_0\approx5260$\,\AA), so lithium
  is not photoelectrically active at 5461\,\AA.
  \item[$b$] The paper states the \emph{observed} $V_{0,\text{obs}} = 0.52$\,V
  with $V_c = 2.51$\,V, giving $V_{0,\text{true}} = 3.03$\,V.
  The computed value 3.07\,V from $h/e \cdot (\nu - \nu_0)$ is consistent
  within the stated 1\% uncertainty for this line.
\end{tablenotes}
\end{table}

\paragraph{Linear regression.}
A linear regression of $V_0$ on $\nu$ using the five most reliable
sodium data points gives:
\begin{itemize}
  \item Slope: $h/e = 4.128\times10^{-15}$\,V$\cdot$s
  \item Intercept: $-W_{\text{Na}}/e = -1.80$\,V (at $\nu = 0$)
  \item Implied $h$ (with modern $e$): $4.128\times10^{-15}\times1.602\times10^{-19}
    = 6.61\times10^{-34}$\,J$\cdot$s (0.5\% high relative to modern value
    due to rounding)
\end{itemize}

\paragraph{Long-wavelength limit and work function.}
Extrapolation of the fitted line to $V_0 = 0$ gives $\nu_0$ (threshold
frequency).  From the contact-EMF-corrected line for sodium,
Millikan found $\nu_0 = 43.9\times10^{13}$\,Hz, corresponding to
$\lambda_0 = 6800$\,\AA\ (deep red).  Direct measurement confirmed
that sodium was photoelectrically active at the yellow mercury lines
$\lambda = 5780, 5790$\,\AA\ and even at the red line $\lambda = 6235$\,\AA,
consistent with $\lambda_0 \approx 6800$\,\AA.

\section{Millikan's Reluctance to Accept the Photon}

A remarkable feature of this episode is that Millikan, while confirming
Einstein's equation to within 0.5\%, refused to accept the physical interpretation.
He wrote in January 1916 (before the definitive result):
\begin{quote}
``Einstein's photoelectric equation \ldots cannot in my judgment be looked upon
at present as resting upon any sort of a satisfactory theoretical foundation.''
\end{quote}
Yet the equation ``actually represents very accurately the behavior.''  In his
1916 paper he describes Einstein's light-quantum hypothesis as ``bold, not to
say the reckless, hypothesis.''

By 1923 (after Compton scattering), Millikan had accepted the photon.
In his 1923 Nobel address he said: ``This work resulted, contrary to my own
expectation, in the first direct experimental proof \ldots of the
Einstein equation'' (quoted by Holton, 1999).

\section{Circuit of Prior Results and Millikan's Innovations}

\begin{table}[ht]
\centering
\caption{Comparison of key pre-Millikan photoelectric measurements of
$h/e$ (or equivalently of $h$ using the known $e$).  Millikan attributes
the systematic underestimation in earlier work to stray short-wavelength
light, restricted frequency range, and uncontrolled contact EMFs.}
\label{tab:comparison}
\begin{tabular}{lllll}
\toprule
Authors & Year & Metal(s) & 
  $h$ obtained ($10^{-27}$\,erg$\cdot$s) & Error \\
\midrule
Hughes               & 1912 & Cu, Zn, Al, Sn, Pb, etc. & 3.55--5.85 & $-$10 to $-$45\% \\
Richardson \& Compton& 1912 & 8 metals incl.\ Na, Al    & $\sim$3.5--5.8 & large \\
Millikan (Na)        & 1916 & Sodium       & 6.569 & $-$0.8\% \\
Millikan (Li)        & 1916 & Lithium      & 6.584 & $-$0.6\% \\
Planck (radiation)   & 1913 & ---          & 6.415 & $-$3.2\% \\
\textbf{Millikan mean} & 1916 & Na + Li    & \textbf{6.57} & $\mathbf{-0.84\%}$ \\
Modern (NIST 2022)   & 2022 & ---          & 6.6261 & 0 (exact) \\
\bottomrule
\end{tabular}
\end{table}

\section{Key References}

\subsection*{Primary Sources}

\begin{itemize}
\item[\textbf{[1]}] P.\ Lenard,
  ``Ueber die lichtelektrische Wirkung,''
  \emph{Annalen der Physik} \textbf{313}(5):149--198 (1902).
  \doi{10.1002/andp.19023130510}
  [Also cited as \emph{Ann.\ Phys.} \textbf{8}:149--198 in the \emph{neue Folge}.]

\item[\textbf{[2]}] A.\ Einstein,
  ``\"Uber einen die Erzeugung und Verwandlung des Lichtes betreffenden
    heuristischen Gesichtspunkt,''
  \emph{Annalen der Physik} \textbf{17}:132--148 (1905).
  \doi{10.1002/andp.19053220607}

\item[\textbf{[3]}] O.\ W.\ Richardson and K.\ T.\ Compton,
  ``The photoelectric effect,''
  \emph{Philosophical Magazine} \textbf{24}:575--594 (1912).
  \doi{10.1080/14786441208637039}

\item[\textbf{[4]}] A.\ Ll.\ Hughes,
  ``On the emission velocities of photo-electrons,''
  \emph{Philosophical Transactions of the Royal Society of London,
    Series A} \textbf{212}:205--226 (1913; submitted 1912).
  \doi{10.1098/rsta.1913.0007}

\item[\textbf{[5]}] R.\ A.\ Millikan,
  ``A Direct Determination of `h',''
  \emph{Physical Review} \textbf{4}(1):73--75 (1914).
  \doi{10.1103/PhysRev.4.73.2}

\item[\textbf{[6]}] R.\ A.\ Millikan,
  ``A Direct Photoelectric Determination of Planck's `h',''
  \emph{Physical Review} \textbf{7}(3):355--388 (1916).
  \doi{10.1103/PhysRev.7.355}
  [Full text: \texttt{archive.org} identifier
   \texttt{sim\_physical-review\_1916-03\_7\_3}, pp.\ 355--388.]
\end{itemize}

\subsection*{Secondary Sources}

\begin{itemize}
\item[\textbf{[7]}] G.\ Holton,
  ``Centennial Focus: Millikan's Measurement of Planck's Constant,''
  \emph{Physics} (APS Physics Focus), vol.\ 3, story 23 (1999).
  \doi{10.1103/physrevfocus.3.23}
  URL: \url{https://physics.aps.org/story/v3/st23}

\item[\textbf{[8]}] B.\ R.\ Wheaton,
  ``Philipp Lenard and the Photoelectric Effect, 1889--1911,''
  \emph{Historical Studies in the Physical Sciences}
  \textbf{9}:299--322 (1978).
  \doi{10.2307/27757381}  JSTOR:27757381.
\end{itemize}

\subsection*{BibTeX Entries}

\begin{verbatim}
@article{Lenard1902,
  author  = {Lenard, Philipp},
  title   = {{\"U}ber die lichtelektrische {W}irkung},
  journal = {Annalen der Physik},
  volume  = {313},
  number  = {5},
  pages   = {149--198},
  year    = {1902},
  doi     = {10.1002/andp.19023130510},
  note    = {Also cited as \emph{Ann.\ Phys.}\ \textbf{8}:149--198
             (neue Folge).  Nobel Prize in Physics, 1905.}
}

@article{Einstein1905,
  author  = {Einstein, Albert},
  title   = {{\"U}ber einen die {E}rzeugung und {V}erwandlung des
             {L}ichtes betreffenden heuristischen {G}esichtspunkt},
  journal = {Annalen der Physik},
  volume  = {17},
  pages   = {132--148},
  year    = {1905},
  doi     = {10.1002/andp.19053220607}
}

@article{RichardsonCompton1912,
  author  = {Richardson, Owen Willans and Compton, Karl Taylor},
  title   = {The photoelectric effect},
  journal = {Philosophical Magazine},
  volume  = {24},
  pages   = {575--594},
  year    = {1912},
  doi     = {10.1080/14786441208637039}
}

@article{Hughes1913,
  author  = {Hughes, Arthur Llewelyn},
  title   = {On the emission velocities of photo-electrons},
  journal = {Philosophical Transactions of the Royal Society
             of London, Series A},
  volume  = {212},
  pages   = {205--226},
  year    = {1913},
  doi     = {10.1098/rsta.1913.0007}
}

@article{Millikan1914,
  author  = {Millikan, Robert Andrews},
  title   = {A Direct Determination of {`$h$'}},
  journal = {Physical Review},
  volume  = {4},
  number  = {1},
  pages   = {73--75},
  year    = {1914},
  doi     = {10.1103/PhysRev.4.73.2}
}

@article{Millikan1916,
  author  = {Millikan, Robert Andrews},
  title   = {A Direct Photoelectric Determination of {Planck's `$h$'}},
  journal = {Physical Review},
  volume  = {7},
  number  = {3},
  pages   = {355--388},
  year    = {1916},
  doi     = {10.1103/PhysRev.7.355},
  note    = {Full text at \texttt{archive.org} identifier
             \texttt{sim\_physical-review\_1916-03\_7\_3}.}
}

@article{Holton1999,
  author  = {Holton, Gerald},
  title   = {Centennial Focus: {M}illikan's Measurement of
             {P}lanck's Constant},
  journal = {Physics (APS Physics Focus)},
  volume  = {3},
  pages   = {story 23},
  year    = {1999},
  doi     = {10.1103/physrevfocus.3.23},
  url     = {https://physics.aps.org/story/v3/st23}
}

@article{Wheaton1978,
  author  = {Wheaton, Bruce R.},
  title   = {Philipp {L}enard and the Photoelectric Effect, 1889--1911},
  journal = {Historical Studies in the Physical Sciences},
  volume  = {9},
  pages   = {299--322},
  year    = {1978},
  doi     = {10.2307/27757381},
  note    = {JSTOR:27757381}
}
\end{verbatim}

\section{Physical Constants Used}

\begin{table}[ht]
\centering
\caption{Physical constants.  Millikan's values (1916) vs.\ modern NIST CODATA
2022.  The 0.65\% error in Millikan's $e$ propagates directly to his $h$.}
\label{tab:constants}
\begin{tabular}{llll}
\toprule
Constant & Millikan (1916) & Modern (NIST 2022) & Error \\
\midrule
$h$ (erg$\cdot$s) & $6.57\times10^{-27}$ & $6.6261\times10^{-27}$ & $-$0.84\% \\
$e$ (esu) & $4.774\times10^{-10}$ & $4.803\times10^{-10}$ & $-$0.60\% \\
$e$ (C) & $1.592\times10^{-19}$ & $1.602\,176\,634\times10^{-19}$ (exact) & $-$0.65\% \\
$h/e$ (V$\cdot$s) & $4.128\times10^{-15}$ & $4.136\times10^{-15}$ & $-$0.19\% \\
\bottomrule
\end{tabular}
\end{table}

\end{document}
